% !TeX program = lualatex
\documentclass[11pt]{article}

\usepackage{iftex}
\usepackage[spanish,es-nodecimaldot]{babel}
\ifPDFTeX
  \usepackage[utf8]{inputenc}
  \usepackage[T1]{fontenc}
  \usepackage{lmodern}
  \usepackage{microtype}
  \newcommand{\essaydisplayfont}{\rmfamily}
\else
  \usepackage{fontspec}
  \defaultfontfeatures{Ligatures=TeX}
  \setmainfont{EB Garamond}[
    UprightFont=EB Garamond Regular,
    ItalicFont=EB Garamond Italic,
    BoldFont=EB Garamond Bold,
    BoldItalicFont=EB Garamond Bold Italic
  ]
  \newfontfamily\essaydisplayfont{EB Garamond}[
    UprightFont=EB Garamond Regular,
    ItalicFont=EB Garamond Italic,
    BoldFont=EB Garamond Bold,
    BoldItalicFont=EB Garamond Bold Italic
  ]
\fi
\usepackage{csquotes}
\usepackage[letterpaper,margin=1in]{geometry}
\usepackage{amsmath,amssymb}
\usepackage{setspace}
\usepackage{caption}
\usepackage{titlesec}
\usepackage{fancyhdr}
\usepackage{xcolor}
\usepackage{hyperref}
\usepackage[backend=biber,style=apa,sorting=nyt,giveninits=true]{biblatex}
\DeclareLanguageMapping{spanish}{spanish-apa}
\addbibresource{biblio.bib}

\definecolor{essayink}{HTML}{1E1A17}
\definecolor{essayaccent}{HTML}{7A5C45}
\definecolor{essayrule}{HTML}{CCB79D}

\setstretch{1.12}
\setlength{\parskip}{0.18em}
\setlength{\parindent}{1.35em}
\setlength{\headheight}{14pt}
\widowpenalty=10000
\clubpenalty=10000
\displaywidowpenalty=10000
\captionsetup{font=footnotesize,labelfont=bf,skip=0.35em}
\hypersetup{colorlinks=true,linkcolor=black,citecolor=black,urlcolor=blue,pdfauthor={OpenAI},pdftitle={Ensayo sobre redes biológicas}}

\title{\textbf{Redes biológicas desde una perspectiva topológica: \\propiedades, modelado y valor analítico en el estudio de los sistemas vivos}}
\author{}
\date{}

\hypersetup{
  colorlinks=true,
  linkcolor=essayink,
  citecolor=essayink,
  urlcolor=essayaccent,
  pdftitle={Redes biológicas}
}
\urlstyle{same}

\titleformat{name=\section,numberless}
  {\essaydisplayfont\Large\bfseries\centering\color{essayink}}
  {}{0pt}{}
\titleformat{name=\subsection,numberless}
  {\essaydisplayfont\large\bfseries\color{essayaccent}}
  {}{0pt}{}
\titlespacing*{\section}{0pt}{1.8em}{0.7em}
\titlespacing*{\subsection}{0pt}{1.2em}{0.45em}

\newcommand{\essaytitle}{Redes biol\'ogicas: propiedades, modelado y valor anal\'itico en el estudio de los sistemas vivos}
\newcommand{\essayrunningtitle}{Redes biol\'ogicas}
\newcommand{\essaytitleblock}{%
  \begin{center}
    {\essaydisplayfont\fontsize{18}{22}\selectfont\bfseries \essaytitle\par}
    \vspace{0.85em}
    %{\color{essayrule}\rule{0.32\linewidth}{0.6pt}\par}
  \end{center}
}

\pagestyle{fancy}
\fancyhf{}
\fancyhead[L]{\footnotesize\itshape \essayrunningtitle}
\fancyhead[R]{\footnotesize\thepage}
\renewcommand{\headrulewidth}{0pt}

\DeclareFieldFormat{doi}{%
  \ifhyperref
    {\url{https://doi.org/#1}}
    {https://doi.org/#1}}
\ExecuteBibliographyOptions{doi=true,url=false,eprint=false}
\renewcommand*{\bibfont}{\normalsize}
\setlength{\bibitemsep}{0.18em}
\AtBeginBibliography{\setstretch{1.12}}
% Ajuste solicitado: poner en cursiva todos los titulos de las obras citadas,
% incluso en tipos de entrada que APA suele dejar en redonda.
\DeclareFieldFormat[article,inbook,incollection,periodical,misc,legislation,legadminmaterial,jurisdiction,inproceedings,legal]{title}{\mkbibemph{#1}\isdot}
\DeclareFieldFormat{citetitle}{\mkbibemph{#1}}

\makeatletter
\DeclareFieldFormat{citehyperref}{%
  \DeclareFieldAlias{bibhyperref}{noformat}%
  \bibhyperref{#1}}
\DeclareFieldFormat{textcitehyperref}{%
  \DeclareFieldAlias{bibhyperref}{noformat}%
  \bibhyperref{%
    #1%
    \ifbool{cbx:parens}
      {\global\boolfalse{cbx:parens}}
      {}}}
\savebibmacro{cite}
\savebibmacro{textcite}
\renewbibmacro*{cite}{%
  \printtext[citehyperref]{%
    \restorebibmacro{cite}%
    \usebibmacro{cite}}}
\renewbibmacro*{textcite}{%
  \ifboolexpr{
    ( not test {\iffieldundef{prenote}}
      and
      ( not test {\ifnumequal{\value{citecount}}{1}}
        or
        not test {\iffieldundef{postnote}} ) )
    or
    test {\ifnameundef{labelname}}
  }
    {\DeclareFieldAlias{textcitehyperref}{noformat}}
    {}%
  \printtext[textcitehyperref]{%
    \restorebibmacro{textcite}%
    \usebibmacro{textcite}}}
\makeatother

\begin{document}

\essaytitleblock

% \section*{Introducción}

Una red biológica representa relaciones entre elementos de un sistema vivo, como proteínas, genes, metabolitos, neuronas, especies o regiones del genoma. En lugar de buscar descripciones individuales, intenta explicar cómo se organizan sus vínculos y la estructura que se forma entre ellos. 

El estudio topológico\footnote{En este ensayo, \enquote{topológico} alude al estudio de la forma sobre las relaciones de la gráfica: quién se conecta con quién, con qué dirección, con qué intensidad y con qué organización global.} permite describir y analizar sistemas mediante nodos, aristas y propiedades globales. 

No se trata solo de aplicar las mismas medidas a cualquier red, sino de interpretar cada propiedad según qué estemos analizando. No significa lo mismo un nodo central en una red de proteínas que en una red trófica o en una red de cromatina.

El análisis de redes biológicas adquiere valor cuando sus propiedades, formas de modelado y aplicaciones se interpretan de manera específica para cada clase de red.\footnote{Aquí la palabra \enquote{clase} se usa en el sentido que designa familias de redes que comparten una misma lógica de representación, aunque cambien la escala, los objetos y el tipo de relaciones.} Por ello, debe preguntarse qué significa cada propiedad dentro de cada clase de red, como las redes de interacción entre proteínas, las redes metabólicas, regulatorias, neuronales, ecológicas o de contacto cromatínico.

% \section*{Argumentación}

% \subsection*{1. Clases de redes}

Bajo esta perspectiva, conviene distinguir varias clases de redes biológicas, cada una definida por una lógica particular de representación. Aunque todas pueden describirse mediante nodos y aristas, difieren en la escala, en el tipo de entidades consideradas y en el significado de sus conexiones. Por ello, la interpretación de sus propiedades no puede separarse de la clase de red a la que pertenecen.

En las redes moleculares se incluyen las redes de interacción entre proteínas, donde una arista representa contacto físico o asociación funcional; las redes metabólicas, donde las conexiones expresan transformaciones bioquímicas entre metabolitos o relaciones entre metabolitos y reacciones; las redes de señalización, que describen cascadas de activación e inhibición; las redes de coexpresión génica, que conectan genes con patrones de expresión semejantes; y las redes de regulación génica, en las que las aristas representan control transcripcional o postranscripcional \parencite{EmmertStreib2014,ZhangHorvath2005,BanfRhee2017,Koutrouli2020}.

En otra escala aparecen las redes neuronales biológicas. En ellas, los nodos pueden ser neuronas individuales o regiones cerebrales, y las aristas pueden representar conectividad anatómica, sinapsis o acoplamiento funcional. Aquí interesa entender cómo el sistema combina segregación local e integración global \parencite{BullmoreSporns2009}.

En una escala mayor se sitúan las redes ecológicas. Entre ellas destacan las redes tróficas, que representan quién consume a quién; las redes de polinización y mutualismo, que describen interacciones entre especies; y las redes sociales animales, útiles para estudiar cooperación, jerarquía, transmisión de información o propagación de patógenos \parencite{Dunne2002,KrauseLusseauJames2009,Bascompte2009}. También pertenecen a esta clase las redes de contacto entre regiones del ADN, que permiten estudiar cómo la proximidad espacial del genoma condiciona la regulación génica \parencite{Beagrie2017}.

Lo importante es que cada clase tiene una interpretación diferente. Una arista puede significar contacto físico, flujo de materia, regulación, correlación, sinapsis, interacción ecológica o proximidad espacial. Precisamente por eso, las mismas medidas topológicas cambian de sentido según la red que se esté estudiando.

% \subsection*{2. Caracterización y propiedades}
Precisamente porque las mismas medidas topológicas cambian de sentido según la red estudiada, es necesario pasar de la clasificación a la caracterización. Esto implica examinar qué propiedades resultan más informativas en cada caso y qué interpretación biológica puede darse a métricas como el grado, la intermediación, el agrupamiento, la modularidad o la longitud de camino. 

En las redes de interacción entre proteínas, un grado alto indica que una proteína participa en muchos complejos o procesos celulares. La intermediación alta nos dice que conecta módulos funcionales distintos y puede actuar como cuello de botella entre procesos. Un agrupamiento alto puede sugerir la presencia de complejos proteicos o ensamblajes moleculares \parencite{Koutrouli2020,Jeong2001,Joy2005}.

Las redes metabólicas se representan mediante redes dirigidas porque una reacción transforma unos compuestos en otros, la relación tiene la orientación de los sustratos hacia los productos. El grado alto de un metabolito puede señalar una posición central en el intercambio bioquímico, pero también puede reflejar metabolitos de moneda metabólica, como ATP, agua o NADH, es decir, moléculas que participan en numerosas reacciones y por ello conectan gran parte de la red. La longitud de camino mide el número mínimo de pasos necesarios para conectar compuestos dentro del espacio de transformaciones bioquímicas; la modularidad aproxima subsistemas bioquímicos; y la existencia de caminos alternativos explica redundancia y plasticidad metabólica \parencite{EmmertStreib2015,Koutrouli2020}.

Las redes de señalización y regulación génica se modelan como redes dirigidas y ponderadas porque importa quién regula a quién. El grado de salida de un factor de transcripción indica el alcance de su influencia regulatoria; el grado de entrada de un gen señala cuántas señales convergen sobre él. Los circuitos de retroalimentación positiva se asocian con memoria y estabilización de estados celulares, mientras que la retroalimentación negativa se vincula con homeostasis y amortiguación. El motivo \textit{feed-forward loop}, es decir, un patrón en el que un regulador controla a un segundo regulador y ambos influyen sobre un mismo gen, se interpreta como un módulo capaz de filtrar ruido o introducir respuestas selectivas en el tiempo. En estas redes, la topología local se interpreta casi siempre en términos dinámicos\footnote{Esto significa que un patrón local de conexiones se describe también por el tipo de comportamiento temporal que tiene, como el tiempo de respuesta, la memoria celular, la estabilidad de estados o transiciones entre estados.} 
\parencite{BanfRhee2017, EmmertStreib2014}. 

% ManganAlon2003,ShenOrr2002,
En las redes de coexpresión una arista no representa regulación directa, sino semejanza estadística entre perfiles de expresión. Por eso, un módulo de coexpresión se interpreta como un conjunto de genes que tienden a expresarse de manera coordinada y que puede asociarse con una condición fisiológica, un tejido, un fenotipo o una enfermedad. Un gen concentrador dentro del módulo puede ser un buen marcador, y la modularidad es más informativa que la centralidad global \parencite{Gnanakkumaar2019,Zheng2021,ZhangHorvath2005}.

En las redes neuronales, el agrupamiento y la modularidad se interpretan como especialización de subredes que procesan tareas distintas. La longitud media de camino y ciertas configuraciones de pequeño mundo se interpretan como eficiencia de integración entre regiones distantes. Los motivos recurrentes pueden señalar microcircuitos de procesamiento, y los nodos altamente conectados pueden comportarse como puntos de alta vulnerabilidad, cuya alteración afecta de manera desproporcionada al conjunto de la red \parencite{Mittal2024,BullmoreSporns2009}. 
% Aquí la propiedad topológica interesa porque se conecta con la organización del procesamiento neuronal, no solo con la forma de la gráfica.


En las redes ecológicas, las propiedades más informativas dependen del subtipo de red. En redes tróficas, la conectancia ayuda a estimar cuán densamente se distribuyen los flujos de energía y materia; en redes de polinización y mutualismo, el anidamiento indica que especies especialistas se relacionan con subconjuntos de especies generalistas; y la modularidad puede mostrar compartimentos ecológicos que limitan o amortiguan perturbaciones \parencite{Dunne2002,Olesen2007,Bascompte2009}. En redes sociales animales, la centralidad de ciertos individuos ayuda a interpretar liderazgo, difusión de información o riesgo de transmisión de patógenos \parencite{KrauseLusseauJames2009}. 
% En este caso, la topología se interpreta directamente en términos de estabilidad, especialización y vulnerabilidad ecológica.

Por último, en las redes de contacto cromatínico, los pesos de las aristas representan frecuencia o intensidad de contacto espacial entre regiones del genoma. Los nodos con muchos contactos pueden señalar regiones o zonas con mucha actividad regulatoria; los módulos reflejan dominios espaciales o compartimentos funcionales; y las interacciones de largo alcance permiten relacionar potenciadores distales con genes objetivo. Estas redes describen cómo la arquitectura tridimensional del genoma condiciona la expresión génica \parencite{Beagrie2017}.
 

% \subsection*{3. Modelado y análisis}


El modelado depende de la clase de red. Las redes de interacción entre proteínas se construyen a partir de métodos experimentales como el ensayo de doble híbrido y la purificación por afinidad acoplada a espectrometría de masas \parencite{Koutrouli2020}.

Las redes metabólicas se derivan de reconstrucciones estequiométricas y pueden representarse como gráficas dirigidas o como redes bipartitas metabolito--reacción \parencite{Koutrouli2020}.

Las redes de señalización y las regulatorias se construyen con transcriptómica, datos de unión de factores de transcripción al ADN, perturbaciones experimentales y conocimiento previo. Estas redes suelen modelarse como redes dirigidas y firmadas, es decir, con aristas que distinguen activación e inhibición. Si el interés está en estados celulares y reglas de activación, son útiles los modelos booleanos; si interesa la cinética, se emplean ecuaciones diferenciales; y si el ruido es relevante, se introducen formulaciones estocásticas para estudiar atractores, estados estables y transiciones entre estados regulatorios \parencite{Elowitz2002,Kauffman1969,BanfRhee2017,EmmertStreib2014}.
% En esta clase, el objetivo no es solo dibujar enlaces, sino reproducir trayectorias, atractores, es decir, 


Las redes de coexpresión se modelan a partir de matrices de expresión. A partir de correlaciones o medidas afines se construyen redes ponderadas, y después se detectan módulos y genes que sean hubs dentro de cada módulo \parencite{ZhangHorvath2005}. 

Las redes neuronales se construyen con trazado anatómico, resonancia magnética de difusión, electrofisiología o series temporales de actividad funcional. Según los datos, las aristas representan fibras anatómicas, sinapsis o acoplamiento funcional entre regiones. El análisis se centra en combinar pesos, distancias, partición modular y comparación entre estados fisiológicos o patológicos \parencite{BullmoreSporns2009}.

Las redes ecológicas se obtienen mediante observación de campo, registros de dieta, visitas florales o seguimiento de interacciones sociales. Según el caso, pueden modelarse como redes dirigidas, bipartitas o ponderadas. Se aplican simulaciones de extinción, análisis de robustez y detección de especies \parencite{Olesen2007,Memmott2004,Bascompte2009}. 

En las redes de cromatina, por último, el modelado parte de matrices de contacto generadas por técnicas especializadas (como Hi-C o GAM), diseñadas para medir contactos entre regiones del genoma, que luego se transforman en gráficas para detectar regiones con interacción frecuente y relaciones de largo alcance \parencite{Beagrie2017}.


% \subsection*{4. Alcance, relevancia y problemas abordados}

El alcance de estas redes depende del estudio. En redes de interacción entre proteínas, el análisis se usa para identificar proteínas esenciales, complejos funcionales y módulos asociados con enfermedades.

En redes metabólicas, para localizar reacciones, rutas alternativas y puntos de fragilidad o plasticidad del metabolismo. 

En redes regulatorias y de señalización, para identificar reguladores, circuitos de respuesta, mecanismos de diferenciación y rutas de desregulación patológica \parencite{BanfRhee2017,EmmertStreib2014,Jeong2001}.

En redes de coexpresión, la relevancia principal está en resumir la variación transcriptómica en módulos interpretables y asociarlos con fenotipos, tejidos o enfermedades. Esto ha sido útil para priorizar biomarcadores y genes candidatos, pero también para distinguir qué parte de la señal es estable y qué parte cambia entre condiciones \parencite{Gnanakkumaar2019,Zheng2021,ZhangHorvath2005}.

En redes neuronales, el análisis de integración, segregación y centralidad ha permitido estudiar cómo se reorganiza el cerebro durante tareas cognitivas, lesiones o trastornos neurológicos \parencite{BullmoreSporns2009}. 

En redes ecológicas, permiten anticipar extinciones secundarias, identificar especies clave, evaluar la estabilidad de comunidades y estimar cómo se propaga una perturbación a través del sistema \parencite{Bascompte2009,Dunne2002,Memmott2004}. 

En redes de cromatina, el objetivo es vincular organización espacial y regulación génica, especialmente cuando se busca explicar por qué un gen se activa o se silencia por contactos a larga distancia \parencite{Beagrie2017}.

Cada propiedad tiene un significado diferente según la clase de red, y cada clase de red permite resolver problemas distintos. Las redes de proteínas y regulación ayudan a priorizar blancos experimentales; las metabólicas permiten estudiar la adaptación bioquímica; las de coexpresión ordenan grandes volúmenes de datos transcriptómicos; las neuronales describen arquitectura funcional y vulnerabilidad cerebral; las ecológicas orientan decisiones de conservación; y las de cromatina conectan la estructura del genoma con control transcripcional. 



% \section*{Conclusión}
Las redes biológicas permiten estudiar sistemas vivos como estructuras relacionales, donde proteínas, genes, metabolitos, neuronas, especies o regiones del genoma adquieren sentido no solo por sus propiedades individuales, sino por la forma en que se conectan entre sí. Consideradas en conjunto, sus distintas clases, propiedades y formas de modelado muestran que el valor del análisis de redes depende de interpretar cada estructura de acuerdo con el sistema biológico que representa.

Su utilidad no está únicamente en calcular métricas, sino en comprender qué significa cada una de ellas dentro de una clase específica de red. Una misma propiedad, como el grado, la centralidad, la modularidad o la longitud de camino, puede tener significados distintos según se aplique a una red de interacción molecular, una red metabólica, una red regulatoria, una red neuronal, una red ecológica o una red de contacto cromatínico. 

No son solo representaciones abstractas de sistemas complejos, sino una herramienta que permite transformar datos en hipótesis interpretables. Al ser correctamente contextualizadas, ayudan a decidir qué medir, qué comparar, qué perturbar y qué priorizar en el laboratorio, en la clínica, en el estudio del cerebro o en la conservación ecológica, lo que permite comprender mejor la dependencia, la vulnerabilidad y la dinámica de los sistemas vivos.

\printbibliography[title={Referencias}]

\end{document}
